\section{分类训练和限时训练}

\subsection{使用说明}

\begin{enumerate}
  \item \textbf{分类训练}：按知识点逐项突破，标注难度（★基础 / ★★中档 / ★★★拔高）
  \item \textbf{限时训练}：严格控制时间，模拟考场节奏
  \item 每题预留空白，建议写完整过程后对照解析（解析见本书配套答案册或线上资源）
  \item 训练完成后再查看答案，标记错题并记录到训练记录表
\end{enumerate}

\subsection{一、三角恒等变换与解三角形}

\begin{enumerate}
  \item[★] 已知 $\sin\alpha = \frac35$，$\alpha\in(\frac{\pi}{2},\pi)$，求 $\cos2\alpha$ 和 $\tan(\alpha+\frac{\pi}{4})$。
  \vspace{1.5cm}
  \item[★] $\triangle ABC$ 中 $a=7$，$b=8$，$\cos B=-\frac17$，求 $\angle A$。
  \vspace{1.5cm}
  \item[★] 已知 $\tan\alpha = 2$，求 $\frac{\sin\alpha+\cos\alpha}{\sin\alpha-\cos\alpha}$ 的值。
  \vspace{1.5cm}
  \item[★] 化简 $\sin15^\circ + \sin75^\circ$。
  \vspace{1cm}
  \item[★★] $f(x)=\sin2x+\sqrt3\cos2x$，求单调递增区间和对称轴方程。
  \vspace{2cm}
  \item[★★] $\triangle ABC$ 中 $2b\cos C = 2a-c$。（1）求 $B$；（2）$b=2\sqrt3$，求面积最大值。
  \vspace{3cm}
  \item[★★] $\triangle ABC$ 中 $\angle A=60^\circ$，$a=2$，求 $\triangle ABC$ 面积的最大值。
  \vspace{2.5cm}
  \item[★★] 已知 $\cos(\alpha-\frac{\pi}{6})+\sin\alpha = \frac{4\sqrt3}{5}$，求 $\sin(\alpha+\frac{7\pi}{6})$。
  \vspace{2cm}
  \item[★★] 在 $\triangle ABC$ 中 $\sin^2A-\sin^2B-\sin^2C = \sin B\sin C$，求 $A$。
  \vspace{2cm}
  \item[★★★] $\triangle ABC$ 中 $a,b,c$ 为角 $A,B,C$ 对边，$S=\frac{\sqrt3}{4}(a^2+b^2-c^2)$，求 $C$；若 $c=2$，求 $S_{\triangle ABC}$ 的最大值。
  \vspace{3cm}
\end{enumerate}

\subsection{二、数列}

\begin{enumerate}
  \item[★] 等差 $\{a_n\}$，$a_1=1$，$a_3=5$，求 $S_{10}$。
  \vspace{1.5cm}
  \item[★] 等比 $\{a_n\}$，$a_1=2$，$a_4=16$，求 $a_n$ 和 $S_5$。
  \vspace{1.5cm}
  \item[★] 已知 $S_n=n^2+n$，求 $a_n$。
  \vspace{1.5cm}
  \item[★★] $a_1=1$，$a_{n+1}=2a_n+1$，求 $a_n$。
  \vspace{2cm}
  \item[★★] $S_n$ 为前 $n$ 项和，$a_1=1$，$S_{n+1}=4a_n+2$，$b_n=a_{n+1}-2a_n$。（1）证 $\{b_n\}$ 等比；（2）求 $a_n$。
  \vspace{3cm}
  \item[★★] 已知 $a_n=\frac1{n(n+1)}$，求 $S_n$ 并证 $S_n<1$。
  \vspace{2cm}
  \item[★★] 已知 $a_n=2n-1$，$b_n=2^{n-1}$，$c_n=a_nb_n$，求 $\{c_n\}$ 前 $n$ 项和。
  \vspace{3cm}
  \item[★★] 等差 $\{a_n\}$ 中 $a_1+a_2+a_3=15$，$a_1a_2a_3=80$，求 $a_n$。
  \vspace{2cm}
  \item[★★★] 已知 $a_1=1$，$a_{n+1}=2a_n+3^n$，求 $a_n$。
  \vspace{2.5cm}
  \item[★★★] $a_1=1$，$a_{n+1}=\frac{n+1}{2n}a_n$，求 $a_n$ 和 $\lim_{n\to\infty} S_n$。
  \vspace{3cm}
\end{enumerate}

\subsection{三、立体几何}

\begin{enumerate}
  \item[★] 正方体 $ABCD-A_1B_1C_1D_1$ 中，求异面直线 $A_1B$ 与 $B_1C$ 所成角。
  \vspace{1.5cm}
  \item[★] 三棱锥 $P-ABC$，$PA\perp$ 底面，$\angle ABC=90^\circ$，$PA=AB=BC=2$。求 $PC$ 与平面 $PAB$ 所成角的正弦值。
  \vspace{3cm}
  \item[★★] 四棱锥 $P-ABCD$，底面矩形，$PA\perp$ 底面，$PA=AD=2$，$AB=3$，$E$ 为 $PD$ 中点。（1）求二面角 $E-AC-B$ 余弦值；（2）求 $B$ 到平面 $AEC$ 距离。
  \vspace{4cm}
  \item[★★] 直三棱柱 $ABC-A_1B_1C_1$，$AB=BC=2$，$AA_1=3$，$\angle ABC=120^\circ$。$D$ 为 $CC_1$ 中点，求二面角 $A-BD-C$ 的余弦值。
  \vspace{4cm}
  \item[★★] 三棱锥 $P-ABC$ 中 $PA=PB=PC=2$，$AB=BC=CA=2\sqrt2$。$M$ 为 $PC$ 中点，求 $BM$ 与平面 $PAC$ 所成角的正弦值。
  \vspace{4cm}
  \item[★★★] 如图，矩形 $ABCD$ 中 $AB=2$，$AD=4$，$E$ 为 $CD$ 中点。$\triangle ADE$ 沿 $AE$ 翻折至 $\triangle AD'E$，$D'$ 在面 $ABCE$ 上的射影 $H$ 在 $BE$ 上。求 $BD'$ 与平面 $AD'E$ 所成角的正弦值。
  \vspace{4cm}
\end{enumerate}

\subsection{四、概率与统计}

\begin{enumerate}
  \item[★] 袋中 3 红 2 白，取 2 个球。求（1）恰有 1 红概率；（2）红球数分布列和期望。
  \vspace{3cm}
  \item[★] 射击 5 次，每次命中 0.8。求（1）恰中 4 次概率；（2）期望和方差。
  \vspace{2.5cm}
  \item[★★] 产品一级品率 0.9，抽 $n$ 件。至少有一件一级品的概率 $\ge0.999$，求 $n_{min}$。
  \vspace{2cm}
  \item[★★] 甲胜率 0.6，乙胜率 0.4，三局两胜制。求甲获胜概率。
  \vspace{2.5cm}
  \item[★★] 某路段汽车速度 $X\sim N(72,6^2)$。（1）求 $P(X>84)$；（2）随机测 3 辆，恰有 1 辆超 84km/h 的概率。
  \vspace{3cm}
  \item[★★] $X$ 的分布列：$P(X=0)=0.2$，$P(X=1)=0.3$，$P(X=2)=0.4$，$P(X=3)=0.1$。求 $E(X)$、$D(X)$、$E(2X+3)$。
  \vspace{2.5cm}
  \item[★★★] 甲袋 3 红 2 白，乙袋 2 红 3 白。先从甲袋取 1 球放入乙袋，再从乙袋取 2 球。求（1）最后取到 1 红 1 白的概率；（2）最后取到红球数的期望。
  \vspace{4cm}
  \item[★★★] 射击运动员命中率 $p$（$0<p<1$），重复射击直到命中 $k$ 次为止。设 $X$ 为总射击次数，求 $X$ 的分布列和期望。
  \vspace{3cm}
\end{enumerate}

\subsection{五、解析几何}

\begin{enumerate}
  \item[★] 椭圆 $\frac{x^2}{4}+\frac{y^2}{3}=1$，$F$ 为右焦点，过 $F$ 的直线交椭圆于 $A,B$，求 $\triangle AOB$ 面积最大值。
  \vspace{3cm}
  \item[★] 求与圆 $x^2+y^2=1$ 相切且过 $(1,0)$ 的直线方程。
  \vspace{2cm}
  \item[★★] 双曲线 $C:\frac{x^2}{a^2}-\frac{y^2}{b^2}=1$ 右焦点 $F(2,0)$，渐近线 $y=\pm\sqrt3x$。（1）求 $C$ 方程；（2）过 $F$ 的直线交 $C$ 于 $A,B$，$M$ 为 $AB$ 中点，$OM\perp AB$ 时求直线 $AB$ 方程。
  \vspace{4cm}
  \item[★★] 抛物线 $y^2=2px$ 焦点 $F$，过 $F$ 的直线交抛物线于 $A,B$，$|AF|=2|BF|$，求直线 $AB$ 的斜率。
  \vspace{3cm}
  \item[★★] 椭圆 $C:\frac{x^2}{4}+y^2=1$，$A(2,0)$，$B(0,1)$。过 $P(0,\frac12)$ 的直线交 $C$ 于 $M,N$，$AM$ 与 $BN$ 交于 $Q$。求 $Q$ 的轨迹方程。
  \vspace{4cm}
  \item[★★] 抛物线 $y^2=4x$ 的焦点为 $F$，$A,B$ 在抛物线上且 $AF\perp BF$，求原点 $O$ 到直线 $AB$ 距离的最小值。
  \vspace{4cm}
  \item[★★★] 椭圆 $\frac{x^2}{4}+\frac{y^2}{3}=1$，$A(-2,0),B(2,0)$。过 $P(1,0)$ 的直线交椭圆于 $M,N$，直线 $AM$ 与 $BN$ 交于 $R$。求证：$R$ 在定直线上。
  \vspace{4cm}
  \item[★★★] 双曲线 $\frac{x^2}{3}-y^2=1$，$F$ 为右焦点。过 $F$ 的直线交双曲线于 $A,B$，$O$ 为原点。若 $\overrightarrow{OA}\cdot\overrightarrow{OB}=3$，求直线 $AB$ 的方程。
  \vspace{4cm}
\end{enumerate}

\subsection{六、函数与导数}

\begin{enumerate}
  \item[★] 求 $f(x)=x^3-3x^2+1$ 的单调区间和极值。
  \vspace{2.5cm}
  \item[★] 求 $f(x)=\frac12x^2-\ln x$ 的最小值。
  \vspace{2.5cm}
  \item[★★] 已知 $f(x)=\ln x - a(x-1)$，讨论 $f(x)$ 的单调性。
  \vspace{3cm}
  \item[★★] $f(x)=e^x-ax$ 在 $[0,1]$ 上单调递减，求 $a$ 的取值范围。
  \vspace{2cm}
  \item[★★] $f(x)=x\ln x$，求 $f(x)$ 在 $[e^{-1},e]$ 上的值域。
  \vspace{2.5cm}
  \item[★★] $f(x)=e^x-ax-1$，若 $f(x)\ge0$ 恒成立，求 $a$ 的取值范围。
  \vspace{3cm}
  \item[★★] 曲线 $y=x^3-3x$ 在 $x=1$ 处的切线与曲线 $y=ax^2$ 相切，求 $a$。
  \vspace{2.5cm}
  \item[★★] 证明：当 $x>0$ 时，$e^x > 1+x+\frac12x^2$。
  \vspace{2.5cm}
  \item[★★★] $f(x)=x^3-3ax^2+3x+1$ 在 $(0,+\infty)$ 上既有极大值又有极小值，求 $a$ 的取值范围。
  \vspace{3cm}
  \item[★★★] $f(x)=\ln x - \frac12ax^2 + x$，$a\in\mathbb{R}$。（1）讨论单调性；（2）若 $f(x)$ 有两个极值点 $x_1,x_2$，求 $a$ 范围；（3）求证：$f(x_1)+f(x_2) < 1$。
  \vspace{5cm}
  \item[★★★] $f(x)=e^x - \ln(x+m)$。（1）$m=0$ 时求极值；（2）$m\le2$ 时证 $f(x)>0$。
  \vspace{5cm}
  \item[★★★] 已知 $a_n=2^n$，$b_n=\frac{a_n}{(a_n-1)(a_{n+1}-1)}$，求 $\sum_{k=1}^n b_k$。
  \vspace{4cm}
  \item[★★★] 已知函数 $f(x)=\ln(1+x) - \frac{ax}{1+x}$ 在 $x=0$ 处取得极值。（1）求 $a$；（2）证明：当 $n\ge2$ 时，$\sum_{k=2}^n \frac1k < \ln n$。
  \vspace{5cm}
\end{enumerate}

\subsection{七、限时训练}

\subsubsection{限时训练一（基础巩固，30 分钟）}

\begin{enumerate}
  \item（5 分）$A=\{x\mid x^2-3x+2\le0\}$，$B=\{x\mid 2^x\ge2\}$，求 $A\cap B$。
  \item（5 分）计算 $\dfrac{1+2i}{3-4i}$。
  \item（5 分）$\triangle ABC$ 中 $a=2$，$b=3$，$\cos C=\frac14$，求 $c$。
  \item（5 分）等差 $\{a_n\}$ 中 $a_2=3$，$a_5=9$，求 $S_8$。
  \item（12 分）$\vec{a}=(1,2)$，$\vec{b}=(-2,1)$。求 $|\vec{a}+\vec{b}|$ 和 $\vec{a}$ 与 $\vec{b}$ 夹角的余弦值。
  \item（12 分）$f(x)=x^3-3x$，求 $f(x)$ 在 $[-2,2]$ 上的最大值和最小值。
\end{enumerate}

\textbf{目标}：30 分钟，40+/44

\subsubsection{限时训练二（中档强化，45 分钟）}

\begin{enumerate}
  \item（5 分）$f(x)=\sqrt{4-x^2}$ 与 $g(x)=\ln(x+1)$ 的定义域交集。
  \item（5 分）$(1+2x)^6$ 展开式中 $x^3$ 的系数。
  \item（5 分）正方体 $ABCD-A_1B_1C_1D_1$ 中 $M$ 为 $CC_1$ 中点，求 $BM$ 与底面所成角的正切值。
  \item（12 分）$f(x)=\sin2x+2\sqrt3\sin^2x$。（1）求最小正周期；（2）求在 $[0,\frac{\pi}{2}]$ 上的值域。
  \item（12 分）$a_1=2$，$a_{n+1}=2a_n-1$。证 $\{a_n-1\}$ 等比，求 $S_n$。
  \item（12 分）甲胜率 0.6，乙胜率 0.4，三局两胜。求甲获胜概率。
\end{enumerate}

\textbf{目标}：45 分钟，40+/51

\subsubsection{限时训练三（综合拔高，60 分钟）}

\begin{enumerate}
  \item（5 分）$f(x)=e^x-ax$ 在 $[0,1]$ 上单调递减，$a$ 的取值范围。
  \item（5 分）过 $A(1,2)$ 与圆 $x^2+y^2=4$ 相切的直线方程。
  \item（5 分）$\overrightarrow{AB}\cdot\overrightarrow{AC}=4$，$|\overrightarrow{AB}|=2$，$|\overrightarrow{AC}|=4$，求 $\triangle ABC$ 面积。
  \item（12 分）椭圆 $C:\frac{x^2}{a^2}+\frac{y^2}{b^2}=1$ 离心率 $e=\frac{\sqrt6}{3}$，过右焦点 $F$ 的直线交 $C$ 于 $A,B$，$|AB|_{\min}=4$。（1）求椭圆方程；（2）若 $OA\perp OB$，求直线 $l$ 方程。
  \item（12 分）$f(x)=\ln x - \frac12x^2 + x$。（1）求单调区间；（2）证 $x>1$ 时 $f(x)<\frac16(x^2-1)$。
\end{enumerate}

\textbf{目标}：60 分钟，30+/39

\subsubsection{限时训练四（综合模拟，90 分钟）}

\textit{本套为完整 22 题模拟卷（150 分，120 分钟）的缩减版，限时 90 分钟完成前 18 题。}

\begin{enumerate}
  \item（5 分）$A=\{x\mid x^2<4\}$，$B=\{y\mid y=2^x\}$，求 $A\cap B$。
  \item（5 分）$z=\frac{2+i}{1-i}$，求 $\bar{z}$。
  \item（5 分）向量 $\vec{a}=(1,m)$，$\vec{b}=(2,3)$，$\vec{a}\perp\vec{b}$，求 $m$。
  \item（5 分）等比 $\{a_n\}$ 中 $a_1=3$，$a_4=24$，求 $S_5$。
  \item（5 分）$(x-\frac{2}{x})^4$ 展开式中常数项。
  \item（5 分）$f(x)=\ln(x^2-2x-3)$ 的单调递增区间。
  \item（5 分）圆 $x^2+y^2-2x+4y-4=0$ 的圆心到直线 $3x-4y+1=0$ 的距离。
  \item（5 分）$X\sim B(10,0.4)$，求 $D(X)$。
  \item（5 分 多选）函数 $f(x)=\sin x\cos x$，则 A. 奇函数 B. 周期 $\pi$ C. 最大值 1 D. 对称轴 $x=\frac{\pi}{4}$
  \item（5 分 多选）下列命题正确的有：A. $\forall x\in\mathbb{R}$，$x^2>0$ B. $\exists x\in\mathbb{R}$，$2^x=x$ C. $a>b\Rightarrow a^2>b^2$ D. $a>b>0\Rightarrow \frac1a<\frac1b$
  \item（5 分）$f(x)=\begin{cases} 2^x-1, & x\le0 \\ f(x-1), & x>0 \end{cases}$，求 $f(\frac52)$。
  \item（5 分）直线 $l:y=kx+1$ 与圆 $x^2+y^2=4$ 相交于 $A,B$，$|AB|=2\sqrt3$，求 $k$。
  \item（10 分）$\triangle ABC$ 中 $a,b,c$ 为角 $A,B,C$ 对边，$\cos A=\frac12$，$b=2$，$c=3$。求 $a$ 和 $\sin B$。
  \item（12 分）$a_1=1$，$a_{n+1}=a_n+2n+1$。（1）求通项；（2）求 $\frac1{a_1}+\frac1{a_2}+\cdots+\frac1{a_n}< \frac54$。
  \item（12 分）四棱锥 $P-ABCD$，底面正方形，$PA\perp$ 底面，$PA=AB=2$，$E$ 为 $PB$ 中点。求 $AE$ 与平面 $PCD$ 所成角的正弦值。
  \item（12 分）甲命中率 0.7，乙命中率 0.6。甲乙各独立射击 1 次。（1）求两人都命中的概率；（2）求至少 1 人命中的概率；（3）设 $X$ 为命中人数，求分布列和期望。
\end{enumerate}

\textbf{目标}：90 分钟，100+/120

\subsubsection{限时训练五（压轴挑战，90 分钟）}

\begin{enumerate}
  \item（5 分）$a=0.3^{0.2}$，$b=0.2^{0.3}$，$c=\log_{0.3}0.2$，比较大小。
  \item（5 分）$f(x)=x^3+ax^2+bx+c$ 在 $x=-2$ 和 $x=1$ 处取极值，$f(-1)=4$，求 $a,b,c$。
  \item（5 分 多选）$f(x)$ 定义域 $\mathbb{R}$，$f(x+1)$ 为奇函数，$f(x+2)$ 为偶函数，则 A. $f(0)=0$ B. $f(x)$ 周期 4 C. $f(2)=0$ D. $f(x)$ 偶函数
  \item（5 分 多选）椭圆 $\frac{x^2}{4}+\frac{y^2}{2}=1$ 的焦点为 $F_1,F_2$，$P$ 为椭圆上动点，则 A. $|PF_1|+|PF_2|=4$ B. $\angle F_1PF_2$ 最大为 $90^\circ$ C. $S_{\triangle PF_1F_2}$ 最大为 $4$ D. $|PF_1||PF_2|$ 最大为 $4$
  \item（5 分）$f(x)=e^x-2x$ 在 $x=0$ 处的切线与曲线 $y=x^2+ax$ 相切，求 $a$。
  \item（5 分）$\tan\alpha+\tan\beta=2$，$\cos(\alpha-\beta)=\frac35$，求 $\cos(\alpha+\beta)$。
  \item（12 分）$\triangle ABC$ 中 $\sin A\sin C = \sin^2B$ 且 $a^2+c^2=3ac$。（1）求 $B$；（2）若 $b=2$，求 $S_{\triangle ABC}$。
  \item（12 分）$S_n$ 为 $\{a_n\}$ 前 $n$ 项和，$a_1=1$，$a_{n+1}=\frac{n+2}{n}S_n$。（1）求 $a_n$；（2）$b_n=\frac1{a_n}$，求 $\sum b_n$ 并证 $<2$。
  \item（12 分）$f(x)=\ln x - ax^2 + x$。（1）讨论单调性；（2）$f(x)$ 有两个极值点 $x_1,x_2$，求 $a$ 范围；（3）求证 $f(x_1)+f(x_2) > 1-2\ln2$。
  \item（12 分）椭圆 $\frac{x^2}{4}+y^2=1$，$A(-2,0),B(2,0)$。过 $P(1,0)$ 的直线交椭圆于 $M,N$，$AM$ 与 $BN$ 交于 $Q$。求 $Q$ 的轨迹方程，并判断是否与椭圆有交点。
\end{enumerate}

\textbf{目标}：90 分钟，60+/78

\subsection{训练记录表}

\begin{center}
\begin{tabular}{c|c|c|c|c}
\textbf{训练类型} & \textbf{正确题数} & \textbf{总题数} & \textbf{得分率} & \textbf{用时} \\ \hline
三角分类 & & & & \\
数列分类 & & & & \\
立体几何分类 & & & & \\
概率统计分类 & & & & \\
解析几何分类 & & & & \\
函数导数分类 & & & & \\
限时训练一 & & & & \\
限时训练二 & & & & \\
限时训练三 & & & & \\
限时训练四 & & & & \\
限时训练五 & & & & \\
\end{tabular}
\end{center}

\textit{每次训练后记录，统计各模块得分率，重点突破薄弱环节。建议每 2 周重做一次错题。}
